%%%%%%%%%%%%%%%%%%%%%%%%%%%%%%%%%
%
%       EXERCÍCIO
%
%%%%%%%%%%%%%%%%%%%%%%%%%%%%%%%%%

\ifdefstring{\atividade}{prova}{%
    \renewcommand{\valorquestao}{\ValorQ\ pontos}
}%

\begin{exercicioBanco}[\valorquestao]
    A função conjunta de probabilidade entre as variáveis aleatórias independentes \(X\) e \(Y\) é apresentada abaixo.
    %
    \[
    \renewcommand{\arraystretch}{1.5}
    \begin{array}{|c|c|c|c|c|}
        \hline
        X\backslash Y & -2 & 0 & 2 & P(X = x) \\
        \hline
        1 &  &  &  & 0,3 \\
        \hline
        2 &  &  & & 0,7 \\
        \hline
        P(Y = y) & 0,2 &  & 0,3 & \\
        \hline
    \end{array}
    \]
    %
    \begin{enumerate}[(a)]
        \item Complete a tabela.
        \item Calcule da esperança e a variância da variável \(2X - Y\).
    \end{enumerate}
\end{exercicioBanco}

